% !TeX program = xelatex
% RAGGit: Single-tenant on-prem RAG system
% Patras Tech Talk, 26 Nov 2026
% Build: latexmk -pdfxe main.tex
\documentclass[aspectratio=169,11pt]{beamer}

\usetheme[titleformat=regular,progressbar=frametitle]{metropolis}
\usepackage{appendixnumberbeamer}
\AtBeginSection{}

% Greek text via XeLaTeX and Windows system fonts.
\usepackage[greek]{babel}
\usepackage{fontspec}
\setmainfont{Segoe UI}
\setsansfont{Segoe UI}
\setmonofont{Consolas}

\usepackage{tikz}
\usetikzlibrary{positioning,arrows.meta,calc}
\usepackage{booktabs}

\setbeamertemplate{page number in head/foot}[totalframenumber]
\setbeamerfont{title}{size=\Large}
\setbeamerfont{subtitle}{size=\normalsize}
\setbeamerfont{author}{size=\normalsize}
\setbeamerfont{date}{size=\small}

% Use the real screenshot when available; otherwise show a clear placeholder.
\newcommand{\screenshot}[3]{% filename, width, height
  \IfFileExists{images/screenshots/#1}{%
    \includegraphics[width=#2]{images/screenshots/#1}%
  }{%
    \begin{tikzpicture}
      \draw[very thick, dashed, mDarkBrown!60, rounded corners=3pt]
        (0,0) rectangle (#2,#3);
      \node[text=mDarkBrown!80, align=center, font=\scriptsize]
        at (0.5*#2,0.5*#3) {ΣΤΙΓΜΙΟΤΥΠΟ\\[2pt]\texttt{images/screenshots/#1}};
    \end{tikzpicture}%
  }%
}

\title{RAGGit}
\author{Ανδρέας Δασκαλόπουλος}
\date{Patras Tech Talk · 26 Νοεμβρίου 2026 · Πάτρα}

\begin{document}

% 20 minutes for the talk. Q&A follows.
\begin{frame}[plain,noframenumbering]
  \vspace*{1cm}
  {\huge\bfseries\color{mDarkTeal}\inserttitle\par}
  \vspace{0.7cm}
  {\Large Single-tenant, on-prem RAG system\par}
  \vspace{1.2cm}
  {\large \insertauthor\par}
  \vspace{0.2cm}
  {\normalsize \insertdate\par}
  \vfill
  \note{1'. Καλησπέρα. Θα σας δείξω το RAGGit, ένα σύστημα που αφήνει μια
  εταιρεία να κάνει ερωτήσεις πάνω στα δικά της έγγραφα. Η βασική ιδέα είναι
  απλή: κρατάμε την αναζήτηση και το μοντέλο μέσα στο ιδιωτικό περιβάλλον της
  εταιρείας και ζητάμε από το σύστημα να δείχνει τις πηγές της απάντησης.
  Θα περάσουμε από τη διαδρομή ενός εγγράφου μέχρι την απάντηση και θα κλείσω
  με όσα μάθαμε φτιάχνοντάς το.}
\end{frame}

%----------------------------------------------------------------------
\section{Το πρόβλημα}

\begin{frame}{Τι σημαίνει «ανεβάζω ένα έγγραφο σε AI»;}
  \begin{itemize}
    \item Το αρχείο και η ερώτηση φεύγουν από τον υπολογιστή και φτάνουν
      στην υπηρεσία του παρόχου.
    \item Ιδιαίτερα για τα αρχεία, αυτά αποθηκεύονται σε vector stores.
    \item Η πολιτική εξαρτάται από την υπηρεσία, το πλάνο και τις ρυθμίσεις:
      εκπαίδευση, διατήρηση, συνεργάτες και συνδεδεμένα εργαλεία.
    \item Παραμένουν και οι κίνδυνοι παραβίασης ή νόμιμου αιτήματος πρόσβασης.
  \end{itemize}
  \note{1'. Σκεφτείτε ένα συνηθισμένο έγγραφο: μια φόρμα προσωπικού, μια
    προσφορά ή ένα συμβόλαιο. Το ανεβάζουμε σε μια υπηρεσία AI και παίρνουμε
    μια γρήγορη απάντηση. Για να λειτουργήσει, το περιεχόμενο πρέπει να
    φτάσει στην υπηρεσία. Από εκεί και πέρα, οι πολιτικές για εκπαίδευση,
    διατήρηση και επεξεργασία από συνεργάτες διαφέρουν ανά πλάνο και ρύθμιση.
    Ένας σύνδεσμος κοινοποίησης ή ένα συνδεδεμένο εργαλείο μπορεί να ανοίξει
    άλλη διαδρομή για τα δεδομένα.\par
    Στην επόμενη διαφάνεια θα δούμε τρία ενδεικτικά παραδείγματα. Οι όροι
    αλλάζουν, γι' αυτό χρειάζεται έλεγχος της επίσημης πολιτικής για τον
    συγκεκριμένο λογαριασμό. Ακόμη και αν η εκπαίδευση είναι απενεργοποιημένη,
    αυτό δεν σημαίνει αυτόματα ότι οι συνομιλίες διαγράφονται ή ότι δεν
    επεξεργάζονται από υποδομές παρόχων.}
\end{frame}

\begin{frame}{Παράδειγμα πολιτικών υπηρεσιών AI}
  \begin{itemize}
    \item OpenAI: στους προσωπικούς λογαριασμούς Free, Plus και Pro, η
      κοινή χρήση για βελτίωση είναι ενεργή από προεπιλογή.
    \item Anthropic: η χρήση συνομιλιών για βελτίωση εξαρτάται από συναίνεση
      και ορισμένες εξαιρέσεις ασφάλειας.
    \item xAI: υπάρχουν ρυθμίσεις εκπαίδευσης. Το Private Chat δεν
      χρησιμοποιείται για εκπαίδευση και διαγράφεται εντός 30 ημερών.
  \end{itemize}
  \note{1'. Τα παραδείγματα ελέγχθηκαν στις 12 Νοεμβρίου 2026 και αφορούν
    συγκεκριμένες πολιτικές καταναλωτικών λογαριασμών. Η OpenAI έχει ενεργή
    από προεπιλογή την κοινή χρήση για βελτίωση, αλλά μπορεί να απενεργοποιηθεί·
    αυτό δεν διαγράφει το ιστορικό, ενώ τα Temporary Chats μπορεί να
    διατηρηθούν έως 30 ημέρες για λόγους ασφάλειας. Η Anthropic περιγράφει
    χρήση μετά από συναίνεση, με ορισμένες εξαιρέσεις· η διατήρηση μπορεί να
    φτάσει τα πέντε χρόνια όταν ο χρήστης επιτρέψει τη χρήση. Στην xAI, το
    Private Chat δεν χρησιμοποιείται για εκπαίδευση και διαγράφεται εντός
    30 ημερών, εκτός αν ισχύουν νομικοί λόγοι, ασφάλεια ή συμμόρφωση.
    Πρόκειται για περιλήψεις, όχι νομική συμβουλή. Ελέγξτε τους όρους για το
    συγκεκριμένο πλάνο και λογαριασμό πριν την ομιλία.}
\end{frame}

%----------------------------------------------------------------------
\section{Η ιδέα}

\begin{frame}{Τι θέλουμε να πετύχουμε}
  \textbf{Στόχος}

  Ένα RAG σύστημα αποκλειστικά για ένα οργανισμό (single-tenant) όπου τα LLM, τα RAG pipelines και τα δεδομένα μας είναι
  όλα τοπικά (on premises).

  \medskip
  \textbf{Προς}

  \begin{itemize}
     \item Υπουργεία, δήμοι, νοσοκομεία, ιατρεία, φορολογικά και δικηγορικά γραφεία.
     \item Μεγάλος όγκος ψηφιακών αρχείων.
     \item Έγγραφα, ερωτήματα και απαντήσεις που πρέπει να παραμένουν ιδιωτικά.
     \item Διατεθειμένοι να καλύψουν το κόστος εξοπλισμού για τα τοπικά AI models.
   \end{itemize}
  \note{1,5'. Το RAGGit απευθύνεται σε οργανισμούς με μεγάλο αρχείο ψηφιακών
    εγγράφων, για παράδειγμα μεγάλα δικηγορικά γραφεία, υπουργεία, δήμους και
    νοσοκομεία. Συχνά αυτά τα έγγραφα περιέχουν πληροφορίες που πρέπει να
    παραμένουν ιδιωτικές. Το ίδιο ισχύει και για τα ερωτήματα και τις
    απαντήσεις. Στόχος είναι να βοηθήσει το προσωπικό να βρίσκει πληροφορίες
    σε αυτά τα αρχεία και να ελέγχει τις πηγές τους. Η τοπική λειτουργία
    απαιτεί workstation και εξοπλισμό για τα AI models και το RAGGit back-end.
    Ο οργανισμός πρέπει να είναι διατεθειμένος να καλύψει αυτό το κόστος.
    Σε λίγο θα δούμε πώς γίνεται αυτή η αναζήτηση.}
\end{frame}

\begin{frame}{Retrieval-Augmented Generation (RAG)}
  \textbf{Ορισμός}

  Μια μέθοδος που δίνει στα μεγάλα γλωσσικά μοντέλα (LLMs) τη δυνατότητα να
  ψάχνουν και να βρίσκουν στοιχεία σε εξωτερικές πηγές δεδομένων, προτού
  διαμορφώσουν την απάντησή τους.

  \medskip
  \textbf{RAG pipelines}

  \begin{itemize}
    \item Data Ingestion (Εισαγωγή των δεδομένων)
    \item Retrieval \& Generation (Ερώτηση και Απάντηση)
  \end{itemize}
  \note{1,5'. Ορισμός με απλά λόγια: αντί να ρωτάμε το μοντέλο μόνο με ό,τι
    θυμάται από την εκπαίδευσή του, του δίνουμε τα σχετικά αποσπάσματα από
    τις δικές μας πηγές και απαντά πάνω σε αυτά. Αυτό είναι το RAG.\par
    Οι RAG pipelines έχουν δύο φάσεις. Η data ingestion ετοιμάζει το υλικό
    μία φορά ανά έγγραφο· η retrieval and generation τρέχει σε κάθε ερώτηση.
    Στις επόμενες δύο διαφάνειες τις βλέπουμε αναλυτικά.}
\end{frame}

\begin{frame}{Data Ingestion}
	\small
  \begin{enumerate}
    \item \textbf{Εξαγωγή Text}\\
	Αφαίρεση περιττών κενών, κεφαλίδων, υποσέλιδων και αριθμών σελίδων.
    \item \textbf{Τεμαχισμός (Chunking)}\\
      Τα μοντέλα διαβάζουν περιορισμένο κείμενο τη φορά, οπότε το σύστημα
      κόβει τα μεγάλα έγγραφα σε κομμάτια (chunks).
    \item \textbf{Μετατροπή σε Διανύσματα (Embeddings)}\\
      Ένα ειδικό μοντέλο (Embedding LLM) μετατρέπει κάθε κομμάτι σε λίστα
      αριθμών (διάνυσμα) που αποτυπώνει το νόημά του.
    \item \textbf{Αποθήκευση στη Διανυσματική Βάση (Vector Database)}\\
      Τα διανύσματα, με το αρχικό κείμενο και τα μεταδεδομένα (τίτλος,
      σελίδα), αποθηκεύονται σε Vector Database.
  \end{enumerate}
  \note{0,5'. Τέσσερα βήματα, μία φορά ανά έγγραφο. Στην πράξη η εξαγωγή
    κειμένου κρύβει εκπλήξεις, ειδικά σε πολύπλοκα PDF.}
\end{frame}

\begin{frame}{Retrieval \& Generation}
  \begin{enumerate}
    \item \textbf{Ανάκτηση (Retrieval)}. Όταν κάνετε μια ερώτηση, το σύστημα τη
      μετατρέπει και αυτήν σε αριθμούς (με το ίδιο Embedding LLM). Ψάχνει στη
      βάση δεδομένων και εντοπίζει τα κομμάτια κειμένου που ταιριάζουν
      μαθηματικά περισσότερο με το νόημα της ερώτησής σας.
    \item \textbf{Επαύξηση (Augmentation)}. Το σύστημα παίρνει την ερώτησή
      σας, κολλάει από κάτω τα σχετικά αποσπάσματα που μόλις βρήκε και τα
      στέλνει μαζί στο Query LLM, το γλωσσικό μοντέλο που συντάσσει την απάντηση.
    \item \textbf{Παραγωγή (Generation)}. Το Query LLM διαβάζει τα συγκεκριμένα
      αποσπάσματα και συντάσσει μια φυσική, τεκμηριωμένη απάντηση, με
      λιγότερες επινοήσεις (hallucinations).
  \end{enumerate}
  \note{1'. Σε κάθε ερώτηση: η ερώτηση γίνεται και αυτή διάνυσμα, με το ίδιο
    embedding μοντέλο. Η βάση επιστρέφει τα κοντινότερα chunks και το query
    LLM συντάσσει την απάντηση μόνο πάνω σε αυτά.}
\end{frame}

\begin{frame}{RAG pipelines}
  \begin{center}
    \begin{tikzpicture}[
      box/.style={rounded corners=2pt, draw=mDarkTeal, fill=mDarkTeal!8,
        font=\scriptsize, align=center, inner sep=4pt, minimum height=9mm},
      hi/.style={box, draw=mLightBrown, fill=mLightBrown!18},
      lab/.style={font=\scriptsize\bfseries, text=mDarkTeal, anchor=west},
      arr/.style={-{Stealth[length=2.2mm]}, semithick, mDarkTeal}
      ]
      \node[lab] at (0,4.0) {Data Ingestion · μία φορά ανά έγγραφο};
      \node[box] (docs) at (1.7,3.1) {Documents\\PDF · DOCX · XLSX};
      \node[box] (text) at (5.0,3.1) {Text extraction\\+ chunking};
      \node[box] (emb) at (8.3,3.1) {Embeddings\\vectors νοήματος};
      \node[hi]  (db) at (11.7,3.1) {Vector database\\LanceDB};
      \draw[arr] (docs) -- (text);
      \draw[arr] (text) -- (emb);
      \draw[arr] (emb) -- (db);

      \node[lab] at (0,1.8) {Query · σε κάθε ερώτηση};
      \node[box] (q) at (1.7,0.9) {Question\\ερώτημα};
      \node[box] (find) at (5.0,0.9) {Retrieval\\επιλογή σχετικών chunks};
      \node[hi]  (llm) at (8.3,0.9) {Generation\\LLM + context};
      \node[hi]  (ans) at (11.7,0.9) {Answer + citations\\απάντηση + παραπομπές};
      \draw[arr] (q) -- (find);
      \draw[arr] (find) -- (llm);
      \draw[arr] (llm) -- (ans);
    \end{tikzpicture}
  \end{center}
  \vspace{2mm}
  \small Η απάντηση βασίζεται στα chunks που επέστρεψε το retrieval και περιλαμβάνει citations προς τις πηγές.
  \note{2,5'. Το RAG έχει δύο φάσεις. Η πρώτη λέγεται ingestion και γίνεται
    όταν προσθέτουμε ένα έγγραφο. Το text extraction βγάζει το κείμενο από το
    αρχείο. Το chunking το χωρίζει σε μικρά κομμάτια, τα chunks. Για κάθε
    chunk φτιάχνουμε ένα embedding, δηλαδή ένα vector που αποτυπώνει το νόημά
    του, και το αποθηκεύουμε στη vector database, εδώ στη LanceDB. Η ingestion
    γίνεται στο παρασκήνιο: το upload μπαίνει σε queue και ο worker ενημερώνει
    την κατάσταση από Queued και Indexing σε Ready ή Failed.\par
    Η δεύτερη φάση γίνεται σε κάθε query. Το retrieval βρίσκει τα chunks που
    ταιριάζουν περισσότερο στην ερώτηση. Το generation στέλνει την ερώτηση
    μαζί με αυτά τα chunks στο LLM, που συντάσσει την απάντηση. Τα citations
    δείχνουν από ποια πηγή προήλθε.\par
    Οι αγγλικοί όροι είναι αυτοί που συναντάμε συνήθως στην ανάπτυξη τέτοιων
    συστημάτων. Το βασικό είναι να θυμόμαστε τη ροή: ingestion, retrieval,
    generation και citations. Αν δεν βρεθεί σχετικό chunk, το σύστημα λέει ότι
    δεν βρήκε σχετική πληροφορία αντί να μαντέψει.}
\end{frame}

%----------------------------------------------------------------------
\section{Η αρχιτεκτονική}

\begin{frame}{Workstation + Desktop clients}
  \begin{center}
    \begin{tikzpicture}[
      emp/.style={rounded corners=2pt, draw=mDarkTeal, fill=mDarkTeal!8,
        font=\scriptsize, align=center, inner sep=4pt, text width=28mm},
      inner/.style={rounded corners=2pt, draw=mDarkTeal!70,
        fill=white, font=\scriptsize, align=center, inner sep=4pt,
        text width=39mm},
      arr/.style={-{Stealth[length=2.4mm]}, very thick, mDarkTeal}
      ]
      \node[emp] (pc1) at (1.5,3.0) {Υπολογιστής εργαζομένου\\Desktop app};
      \node[emp] (pc2) at (1.5,1.0) {Υπολογιστής εργαζομένου\\Desktop app};
      \draw[very thick, rounded corners=4pt, fill=mLightBrown!12,
        draw=mLightBrown!80] (4.2,0.1) rectangle (13.3,4.1);
      \node[font=\small\bfseries, text=mDarkBrown] at (8.75,3.8)
        {ON-PREM AI WORKSTATION};
      \node[inner] (api) at (6.35,2.75) {REST APIs\\RAG pipelines};
      \node[inner] (model) at (10.95,2.75) {Ollama\\τοπικά query και embedding LLM};
      \node[inner] (vec) at (6.35,1.0) {LanceDB\\διανύσματα αναζήτησης};
      \node[inner] (sql) at (10.95,1.0) {SQLite + αρχεία\\δεδομένα εφαρμογής};
      \draw[arr] (pc1.east) -- node[above,font=\tiny] {LAN} (4.2,3.0);
      \draw[arr] (pc2.east) -- (4.2,1.0);
      \draw[arr,thin] (api) -- (model);
      \draw[arr,thin] (api) -- (vec);
      \draw[arr,thin] (api) -- (sql);
    \end{tikzpicture}
  \end{center}
  \vspace{1mm}
  \note{3'. Αυτό είναι το βασικό διάγραμμα του συστήματος. Ο εργαζόμενος
    χρησιμοποιεί μια εφαρμογή Windows. Η εφαρμογή είναι λεπτή: εμφανίζει τις
    οθόνες και επικοινωνεί με το API μέσω HTTPS. Δεν κατεβάζει τα έγγραφα για
    να τρέξει το μοντέλο στον υπολογιστή του.\par
    Το API, το μοντέλο Ollama, η βάση LanceDB και η βάση SQLite λειτουργούν
    μέσα στην ίδια εγκατάσταση. Τα έγγραφα αποθηκεύονται στο storage της
    εγκατάστασης. Η LanceDB κρατά τα δεδομένα που χρειάζεται η αναζήτηση και
    η SQLite κρατά πληροφορίες για έγγραφα, χρήστες και ερωτήσεις.\par
    Στην παρουσίαση εστιάζουμε σε on-prem εγκατάσταση. Το workstation
    βρίσκεται φυσικά στον χώρο του οργανισμού και συνδέεται στο εταιρικό LAN.
    Δεν χρησιμοποιείται VM σε ιδιωτικό cloud. Το αίτημα πηγαίνει από τον
    client στο API του workstation, και η απάντηση παράγεται εκεί. Έτσι οι
    υπολογιστές των εργαζομένων χρειάζονται μόνο το desktop app και σύνδεση
    στο LAN. Το workstation συγκεντρώνει το μοντέλο και τα δεδομένα, γι' αυτό
    η διαθεσιμότητα και η ασφάλειά του είναι ευθύνη του οργανισμού.}
\end{frame}

\begin{frame}{LLM: τοπικά μοντέλα μέσω Ollama}
  \begin{columns}[T]
    \begin{column}{0.48\textwidth}
      \begin{block}{Embedding LLM}
        \small
        \begin{itemize}
          \item BGE-M3 (1024 διαστάσεις)
          \item Μετατρέπει chunks και ερωτήσεις σε διανύσματα
          \item Το ίδιο μοντέλο και στις δύο φάσεις
        \end{itemize}
      \end{block}
    \end{column}
    \begin{column}{0.48\textwidth}
      \begin{block}{Query LLM}
        \small
        \begin{itemize}
          \item Qwen3 4B Instruct
          \item Απαντά μόνο με βάση τα αποσπάσματα
        \end{itemize}
      \end{block}
    \end{column}
  \end{columns}
  \vspace{2mm}
  \small Τα μοντέλα είναι \textbf{configurable} (στους \texttt{appsettings}
  του workstation): σε ισχυρότερο hardware επιλέγεται μεγαλύτερο Query LLM,
  χωρίς αλλαγή στον κώδικα. Τα δικά μας είναι μικρά, γιατί η ανάπτυξη έγινε
  σε έναν Intel NUC «Skull Canyon» — οι επιλογές περιορίστηκαν από το υλικό
  που είχαμε διαθέσιμο.
  \vspace{1mm}

  \small Όλα εκτελούνται τοπικά μέσω Ollama — καμία κλήση σε εξωτερική υπηρεσία.
  \note{2'. Χρησιμοποιούμε δύο μοντέλα. Το embedding μοντέλο BGE-M3 παράγει
    διανύσματα 1024 διαστάσεων και χρησιμοποιείται και στην εισαγωγή των
    εγγράφων και στην κωδικοποίηση της ερώτησης — πρέπει να είναι το ίδιο
    αλλιώς η σύγκριση δεν έχει νόημα.\par
    Το query μοντέλο είναι το Qwen3 4B Instruct. Το επιλέξαμε μετά από
    μέτρηση στο δικό μας υλικό: τα μικρότερα μοντέλα ήταν γρήγορα αλλά
    έκαναν συγκεκριμένα λάθη στα ελληνικά. Το Qwen3 ήταν το μόνο που δεν
    δήλωσε ψευδές γεγονός. Το τίμημα είναι ο χρόνος στις ευρείες ερωτήσεις.\par
    Αξίζει να πούμε γιατί τόσο μικρά μοντέλα: όλη η ανάπτυξη έγινε σε έναν
    Intel NUC «Skull Canyon» — ένα μικροσκοπικό μηχάνημα με περιορισμένη GPU
    (4 πυρήνες, με τα 11 από τα 37 layers του μοντέλου να φτάνουν στην
    GPU). Σε τέτοιο υλικό δεν χωράνε μεγάλα μοντέλα, οπότε δοκιμάσαμε ό,τι
    χωρούσε και μετρήσαμε.\par
    Σημαντικό: τα μοντέλα δεν είναι hard-coded. Επιλέγονται από τη ρύθμιση
    του workstation (Ollama section στους appsettings) — Embedding LLM και
    Query LLM είναι ονόματα μοντέλων στη configuration, όχι στον κώδικα. Με
    καλύτερο workstation, αλλάζεις δύο γραμμές ρυθμίσεων και ανεβάζεις
    μοντέλο, χωρίς rebuild ή αναφορά σε υπηρεσία του κατασκευαστή. Το μόνο
    που χρειάζεται προσοχή είναι η διάσταση του embedding να ταιριάζει με
    τη βάση (384, 768 ή 1024 ανάλογα με το μοντέλο).\par
    Και τα δύο τρέχουν τοπικά μέσω Ollama. Καμία ερώτηση δεν φεύγει από τον
    οργανισμό.}
\end{frame}

\begin{frame}{Βάσεις δεδομένων}
  \begin{columns}[T]
    \begin{column}{0.48\textwidth}
      \begin{block}{LanceDB — διανύσματα}
        \small
        \begin{itemize}
          \item Embeddings + κείμενο chunks + metadata
          \item ANN αναζήτηση (cosine, HNSW index)
          \item Τοπική αποθήκευση σε αρχεία στον δίσκο
        \end{itemize}
      \end{block}
    \end{column}
    \begin{column}{0.48\textwidth}
      \begin{block}{SQLite — δεδομένα εφαρμογής}
        \small
        \begin{itemize}
          \item Έγγραφα, καταστάσεις ingest, χρήστες
          \item Chunks (παιδιά/γονείς) και ιστορικό ερωτήσεων
          \item Ένα αρχείο \texttt{rag.db}, χωρίς server
        \end{itemize}
      \end{block}
    \end{column}
  \end{columns}
  \vspace{2mm}
  \small Τα πρωτότυπα αρχεία των εγγράφων φυλάσσονται στο ίδιο workstation
  (ένα αρχείο ανά document id).
  \note{2'. Δύο βάσεις με διαφορετικούς ρόλους. Η LanceDB είναι η
    διανυσματική βάση: κρατά τα embeddings μαζί με το κείμενο του chunk και
    μεταδεδομένα, και κάνει την αναζήτηση πλησιέστερων γειτόνων με cosine
    ομοιότητα. Όταν η βιβλιοθήκη μεγαλώσει, χτίζει HNSW index για ταχύτητα.\par
    Η SQLite κρατά ό,τι είναι σχεσιακό: έγγραφα και κατάστασή τους, χρήστες
    και ρόλους, τα chunks με τη σχέση παιδιού-γονιού και το ιστορικό
    ερωτήσεων. Είναι ένα απλό αρχείο, δεν χρειάζεται ξεχωριστός server.\par
    Τα πρωτότυπα uploads μένουν ως ξεχωριστά αρχεία στον δίσκο, ένα για κάθε
    document id, ώστε ο χρήστης να μπορεί να ανοίξει την πρωτότυπη πηγή.}
\end{frame}

\begin{frame}{RAGGit desktop app}
  \begin{center}
    \screenshot{ask-citations.png}{0.72\textwidth}{5.6cm}
  \end{center}
  \note{2,5'. Εδώ φαίνεται η οθόνη Ask του desktop app. Ο εργαζόμενος γράφει
    την ερώτηση στα ελληνικά ή σε άλλη γλώσσα. Η εφαρμογή δείχνει τα βήματα
    της αναζήτησης: πρώτα καταλαβαίνει το αίτημα, μετά ψάχνει, εμφανίζει τις
    πηγές και συντάσσει την απάντηση.\par
    Οι παραπομπές οδηγούν στα αποσπάσματα που χρησιμοποιήθηκαν. Έτσι ο
    εργαζόμενος μπορεί να ελέγξει αν η απάντηση στηρίζεται πράγματι στο
    έγγραφο. Αν δεν υπάρχει σχετικό περιεχόμενο, το σύστημα πρέπει να πει
    ακριβώς ότι δεν βρήκε σχετικό περιεχόμενο — \texttt{no relevant content
    found} — και όχι να συμπληρώσει το κενό με μια πιθανή απάντηση. Δείχνω
    πώς η παραπομπή ανοίγει το σχετικό απόσπασμα και πώς ο χρήστης μπορεί να
    συγκρίνει την απάντηση με την πηγή.\par
    Οι παραπομπές βοηθούν στον έλεγχο, αλλά δεν εγγυώνται από μόνες τους ότι
    κάθε απάντηση είναι σωστή. Ο χρήστης πρέπει να ανοίγει την πηγή όταν η
    πληροφορία έχει σημασία.}
\end{frame}

\begin{frame}{Technology stack}
  \begin{columns}[T]
    \begin{column}{0.48\textwidth}
      \begin{block}{Εφαρμογή}
        \small
        \begin{itemize}
          \item C\# / .NET 10 — ASP.NET Core Web API
          \item WPF + WPF-UI (Fluent, MVVM) client
          \item PdfPig: εξαγωγή κειμένου από PDF
          \item JWT αυθεντικοποίηση με ρόλους
        \end{itemize}
      \end{block}
    \end{column}
    \begin{column}{0.48\textwidth}
      \begin{block}{Δοκιμές και παράδοση}
        \small
        \begin{itemize}
          \item TDD: 700+ τεστ σε τέσσερα επίπεδα
          \item GitHub Actions CI σε κάθε αλλαγή
          \item Spec-driven ανάπτυξη ανά feature
          \item Docker Compose με TLS
        \end{itemize}
      \end{block}
    \end{column}
  \end{columns}
  \vspace{2mm}
  \small Ο κώδικας είναι ανοιχτός:
  \href{https://github.com/daskas-welt/raggit}{github.com/daskas-welt/raggit}
  \note{1'. Όλο το σύστημα είναι γραμμένο σε C\# πάνω στο .NET 10: το API του
    workstation είναι ένα ASP.NET Core project και ο client μια WPF εφαρμογή
    με το WPF-UI kit για Fluent εμφάνιση και το MVVM pattern για τον
    διαχωρισμό της λογικής από την εμφάνιση. Το PdfPig κάνει την εξαγωγή
    κειμένου από τα PDF — είδαμε νωρίτερα ότι αυτή η φάση κρύβει
    εκπλήξεις με πολυτονικά και πολυστήλιχα έγγραφα. Η πρόσβαση γίνεται με
    JWT tokens και ρόλους, ώστε κάθε χρήστης βλέπει μόνο τα δικά του
    έγγραφα.\par
    Στη δεξιά στήλη είναι η διαδικασία: γράφουμε τα τεστ πριν από τον κώδικα
    — πάνω από 700 τεστ σε τέσσερα επίπεδα — και το CI τρέχει σε κάθε
    αλλαγή. Κάθε feature ξεκινά από μια προδιαγραφή που γίνεται plan και
    tasks, ώστε οι αποφάσεις να μένουν καταγεγραμμένες. Το deployment είναι
    ένα Docker Compose stack με TLS.\par
    Όλος ο κώδικας είναι δημόσιος στο GitHub — ο σύνδεσμος είναι και στην
    τελευταία διαφάνεια.}
\end{frame}

\begin{frame}{Τι διαφοροποιεί το RAGGit}
  \begin{columns}[T]
    \begin{column}{0.48\textwidth}
      \begin{block}{Single-tenant, on-prem}
        \small
        \begin{itemize}
          \item Μια εγκατάσταση ανά οργανισμό
          \item Δεδομένα και μοντέλα στον χώρο του πελάτη
          \item Έλεγχος πρόσβασης και αντιγράφων ασφαλείας
        \end{itemize}
      \end{block}
    \end{column}
    \begin{column}{0.48\textwidth}
      \begin{block}{Συντονισμός στον πελάτη}
        \small
        \begin{itemize}
          \item Query / embedding LLM — επιλογή με μέτρηση
          \item Chunk size και TopK ανά βιβλιοθήκη
          \item Δοκιμή στο δικό του υλικό
          \item Όλα ρυθμίσεις, όχι κώδικας
        \end{itemize}
      \end{block}
    \end{column}
  \end{columns}
  \vspace{2mm}
  \small Η αξία: \textbf{μετρημένη ρύθμιση} στα δεδομένα και το υλικό του
  πελάτη — δουλειά σε \textbf{συμβουλευτικό επίπεδο}.
  \note{1'. Το διαφοροποιητικό δεν είναι κάποιος αλγόριθμος· είναι ό,τι όλα
    τα σημεία που συνήθως μαντεύονται εδώ γίνονται μετρήσεις. Το query μοντέλο
    το επιλέξαμε μετρώντας υποψήφια μοντέλα σε πραγματικό ελληνικό corpus —
    τα μικρότερα ήταν γρήγορα αλλά έλεγαν ψευδή γεγονότα με σιγουριά. Το TopK
    το ίδιο: σε μικρή βιβλιοθήκη αρκεί το 3, σε μεγάλη το 5 κρατά την
    ανακτήσιμη απάντηση. Ο chunking και το embedding μοντέλο ακολουθούν την
    ίδια λογική. Και όλα είναι appsettings, όχι αλλαγές στον κώδικα.\par
    Σε κάθε πελάτη η ίδια διαδικασία: ανεβάζουμε τα δικά του έγγραφα, μετράμε
    στο δικό του υλικό, διαλέγουμε μοντέλα και chunking, και παραδίδουμε μια
    εγκατάσταση συντονισμένη στα δικά του δεδομένα. Γι' αυτό μιλάμε για
    δουλειά σε συμβουλευτικό επίπεδο, όχι για απλή πώληση λογισμικού.}
\end{frame}

\begin{frame}{Το κρατάμε ως έχει; Σενάρια και εναλλακτικές}
  \small
  \textbf{Κρατάμε το RAGGit, αλλάζουμε μόνο το deployment}
  \begin{itemize}
    \item Ίδια εγκατάσταση on-prem, όπως παρουσιάστηκε.
    \item Private cloud: VM με GPU που ελέγχει ο οργανισμός — ήδη
      υποστηρίζεται με ρύθμιση.
    \item GPU VM παρόχοι:
      \href{https://vast.ai/}{Vast.AI},
      \href{https://www.paperspace.com/machines}{DigitalOcean Paperspace},
      \href{https://learn.microsoft.com/en-us/azure/architecture/guide/compute/high-performance-computing}{Azure HPC}
  \end{itemize}
  \smallskip
  \textbf{Αντικαθιστούμε το RAGGit με έτοιμη λύση}
  \begin{itemize}
    \item \href{https://ollama.com/}{Ollama} ή
      \href{https://lmstudio.ai/}{LM Studio} μόνο: τοπικό μοντέλο χωρίς
      βιβλιοθήκη εγγράφων.
    \item \href{https://azure.microsoft.com/products/ai-services/ai-foundry}{Azure OpenAI}:
      ισχύς και πολυπλοκότητα, αλλά τα δεδομένα φεύγουν από τον οργανισμό.
    \item \href{https://github.com/infiniflow/ragflow}{RagFlow} (open
      source) και \href{https://gendox.dev/}{Gendox}: έτοιμες RAG πλατφόρμες.
  \end{itemize}
  \note{1,5'. Υποθετική συζήτηση: δείχνω ότι έχω σκεφτεί τα εναλλακτικά
    σενάρια, όχι ότι περιμένω αυτές τις ερωτήσεις. Κρατώ τις απαντήσεις
    σύντομες.\par
    Κρατάμε ως έχει (on-prem): μέγιστος έλεγχος — δικό μας μηχάνημα μέσα στο
    γραφείο. Και δεν χρειάζεται μεγάλο μηχάνημα: η ανάπτυξη έγινε σε Intel
    NUC. Ξεκινάς με ό,τι έχεις και αναβαθμίζεις· τα μοντέλα είναι ρύθμιση,
    όχι κώδικας.\par
    Private cloud: ο δρόμος του μέσου όρου — δικά σου δεδομένα σε υποδομή
    που ελέγχεις, αλλά εκτός γραφείου. Το RAGGit το υποστηρίζει ήδη ως
    δεύτερο mode του deployment: ίδιο artifact set, μόνο αλλαγές σε
    configuration (BIND_ADDRESS, TLS, tailnet), όχι νέος κώδικας. Όποιος
    δεν θέλει να αγοράσει μηχάνημα παίρνει το ίδιο προϊόν.\par
    Vast.AI / Paperspace / Azure HPC: πού νοικιάζεις το GPU VM για το
    σενάριο private cloud. Ίδια λογική, μηνιαίο κόστος αντί για κεφάλαιο.\par
    Ollama / LM Studio: δίνουν το τοπικό μοντέλο, αλλά όχι RAG — δεν έχουν
    βιβλιοθήκη εγγράφων, chunking, ανάκτηση, παραπομπές, χρήστες και ρόλους.
    Το RAGGit χτίζει όλη αυτή τη στρώση πάνω σε Ollama.\par
    Azure OpenAI: μοντέλα OpenAI σε cloud της Microsoft, με υποσχέση ότι δεν
    χρησιμοποιούνται για εκπαίδευση. Τα δεδομένα όμως φεύγουν από τον
    οργανισμό — το ακριβό σημείο της ομιλίας.\par
    RagFlow: ανοιχτού κώδικα RAG engine με έτοιμο interface — σύγκριση με
    το RAGGit ως προς τον έλεγχο, το single-tenant και το βάρος της
    εγκατάστασης. Gendox: εμπορική εναλλακτική RAG-as-a-service — σύγκριση
    κόστους και ελέγχου.\par
    Κλείσιμο: αν ο οργανισμός θέλει πλήρη έλεγχο και μπορεί να έχει ένα
    μηχάνημα, κρατάμε on-prem. Αν όχι, private cloud με την ίδια
    εγκατάσταση. Έτοιμη πλατφόρμα μόνο όταν ο πλήρης έλεγχος δεν μετράει.}
\end{frame}

\begin{frame}[plain,noframenumbering]
  \vfill
  \centering
  {\Large\bfseries\color{mDarkTeal} Η αναζήτηση γίνεται στο εταιρικό περιβάλλον.\par}
  \vspace{5mm}
  {\normalsize Τα έγγραφα και το μοντέλο μένουν στην εγκατάσταση.\par
    Η εταιρεία διαχειρίζεται την πρόσβαση και τα αντίγραφα ασφαλείας.\par}
  \vspace{8mm}
  {\Large\bfseries Ευχαριστώ για την προσοχή!\par}
  \vspace{4mm}
  {\Large\bfseries Ερωτήσεις;\par}
  \vfill
  {\footnotesize Ανδρέας Δασκαλόπουλος · Patras Tech Talk · 26 Νοεμβρίου 2026\par}
  \vspace{2mm}
  {\footnotesize \href{https://github.com/daskas-welt/raggit}{github.com/daskas-welt/raggit}}
  \note{0,5'. Κλείνω με το βασικό σημείο: η ιδιωτικότητα δεν προκύπτει μόνο
    από το μοντέλο. Προκύπτει από το πού τρέχει, ποιος έχει πρόσβαση, πώς
    αποθηκεύονται τα αρχεία και πώς ελέγχονται οι απαντήσεις. Το RAGGit
    μεταφέρει αυτόν τον έλεγχο στην εταιρεία. Ευχαριστώ όλους όσοι ήρθαν
    απόψε και όσους θα δουν την παρουσίαση αργότερα. Από εδώ και πέρα,
    ευχαρίστως να απαντήσω στις ερωτήσεις σας.}
\end{frame}

\end{document}
